\section{Introduction}
\label{sec:introduction}

Modern natural language processing (NLP) is dominated by large pretrained models
such as GPT and BERT, which abstract away the mathematical machinery underneath.
This project takes the opposite approach: every model is built from first principles
in PyTorch, with no pretrained weights, no high-level NLP frameworks, and no
external text corpora beyond what is generated programmatically.

The project is organized into two self-contained tasks:

\begin{itemize}
    \item \textbf{Task 1 — Text Generation.} Six models are trained to predict and
          generate free-form text, covering three architectures (Markov Chain, RNN,
          LSTM) each applied at two granularities (character-level and word-level).
          This progression illustrates how increasingly expressive models capture
          longer-range linguistic dependencies.

    \item \textbf{Task 2 — Chatbots.} Two sequence-to-sequence models are trained
          on a question-answering dataset: an LSTM encoder-decoder with Bahdanau
          attention, and a full Transformer encoder-decoder with multi-head
          self-attention and sinusoidal positional encoding.
\end{itemize}

The goals of this project are to:
\begin{enumerate}
    \item Understand the fundamental trade-offs between probabilistic and
          neural approaches to language modelling.
    \item Observe the practical effect of the vanishing gradient problem, and
          how LSTM gating mechanisms address it.
    \item Compare the convergence behaviour of recurrent and attention-based
          architectures under limited data conditions.
    \item Produce a reproducible, end-to-end pipeline that can be re-run with
          a single command (\texttt{python main.py}).
\end{enumerate}
